\documentclass[11pt]{article}

\usepackage[T1]{fontenc}
\usepackage{amsmath}
\usepackage{amssymb}
\usepackage{booktabs}

\title{What a Standard Error Assumes}
\author{Example Author}
\date{}

\begin{document}
\maketitle

A standard error is not a property of an estimate. It is a property of a model of how the estimate would vary, and that model is chosen, rarely stated, and almost never checked. The number in parentheses under a coefficient is the output of a second model sitting quietly behind the first.

This note is written in \LaTeX, and LibrePaper stored the source as it was published: nothing compiled it on the way. The compiler runs in the browser instead, so the PDF beside the source is the document you are reading.

\section{The second model}

Write the regression as $y = X\beta + \varepsilon$ and the estimator as $\hat{\beta} = (X'X)^{-1}X'y$. Everything about the point estimate is settled by that line. Nothing about its uncertainty is, because the variance of $\hat{\beta}$ is
\begin{equation}
  \operatorname{Var}(\hat{\beta}) = (X'X)^{-1} X' \Omega X (X'X)^{-1},
  \label{eq:sandwich}
\end{equation}
and $\Omega = \operatorname{Var}(\varepsilon)$ is a matrix the data never show you. Every standard error you have ever read is equation~\eqref{eq:sandwich} with some guess plugged in for $\Omega$. The guess is the second model.

The classical choice is $\Omega = \sigma^2 I$: every error has the same variance and none of them are correlated. Under it the sandwich collapses to $\sigma^2 (X'X)^{-1}$, which is the formula in the textbook. It is also the strongest assumption on the menu, and the one most software makes by default.

\section{Three guesses}

The table below is the same coefficient three times. Nothing about the data or the point estimate changes between rows; only the guess about $\Omega$ does.

\begin{table}[h]
  \centering
  \begin{tabular}{lrl}
    \toprule
    Guess about $\Omega$ & Std.\ error & What it assumes \\
    \midrule
    $\sigma^2 I$ & 0.041 & Same variance, no correlation \\
    $\operatorname{diag}(\hat{\varepsilon}_i^2)$ & 0.058 & Any variances, no correlation \\
    Block diagonal by cluster & 0.117 & Anything within, nothing across \\
    \bottomrule
  \end{tabular}
  \caption{One estimate, three standard errors. The estimate is $0.09$ in every row.}
  \label{tab:three}
\end{table}

The third row is nearly three times the first. That is not a small correction to a known number. It is the difference between a result and a shrug, and it came from a decision that most papers describe in half a sentence, if at all.

\section{What clustering assumes}

The cluster-robust estimator replaces $\Omega$ with a block-diagonal matrix built from the residuals within each group. It is honest about correlation inside a cluster and silent about correlation between clusters. It also treats the clusters as the units that were sampled, which is the assumption to read twice.

Two things follow. First, the standard error is only as good as the number of clusters, not the number of observations: forty states are forty draws, however many people live in them. Second, the level of clustering is a claim about the design, not a robustness check. Clustering at a finer level than the one the treatment varies at does not make the interval conservative. It makes it wrong, and usually too narrow.

\section{Reading the parentheses}

When a table reports a standard error, ask what is in $\Omega$ before asking whether the interval excludes zero. The three questions worth asking are short.

\begin{enumerate}
  \item Which rows of the data are allowed to be correlated, and is that the level the treatment varies at?
  \item How many independent units does that leave? Fewer than fifty is a small sample whatever the row count says.
  \item Would the conclusion survive the next row of Table~\ref{tab:three}?
\end{enumerate}

A standard error that survives all three is a claim. One that survives none of them is a number that was printed because the software prints one.

\end{document}
